\chapter[Für mich]{„Für mich“}
\quellseite{9}

Die „Marieliese“, ein schönes, neues Segelschiff, befand sich auf ihrer ersten Reise von Bremen nach Valparaiso in Chili. Alles war bisher günstig verlaufen und fast die Hälfte der Fahrt zurückgelegt. Aber eine dunkle Wolke lagerte über der „Marieliese“: Peter Jensen, der Capitän, lag seit einigen Tagen schwer krank in seiner Koje, und es schien, als sollte er den Hafen nicht erreichen. Aechzend und stöhnend warf er sich im hohen Fieber hin und her oder lag dazwischen hinein bald besinnungslos da. Etwas schien den Kranken zu quälen und zu verfolgen. Eines Tages sagte er zu dem neben ihm stehenden ersten Steuermann, der an seiner Stelle das Commando hatte übernehmen müssen: „Mit mir geht’s zu Ende; ich weiß, ich werde den Hafen nicht mehr erreichen. Aber wie komme ich wohl zum andern Hafen — droben? Ach, helfen Sie mir doch!“ Der Kranke richtete sich jäh auf und packte mit angsterfülltem Blick den Arm des Steuermanns: „Helfen Sie mir, sagen Sie mir, wie ich es machen soll, daß ich nicht verloren gehe.“ „Capitän“, erwiderte der Steuermann, „das weiß ich selbst nicht, habe selbst nie daran gedacht. Ich habe immer rechtschaffen gelebt (?) und meine Pflicht gethan (?) und mich nicht um Gott oder göttliche Dinge bekümmert.“ „So schicken Sie mir den zweiten Steuermann“, befahl der Capitän. Der zweite Steuermann kam, aber auch er konnte ebensowenig wie der erste die Seelenangst des Capitäns beschwichtigen. Ein Mann der Besatzung um den andern wurde gerufen, aber keiner konnte dem Sterbenden den Weg des Heils weisen. Nur der Schiffsjunge war noch nicht in der Kajüte des Capitäns gewesen, aber auch ihn ließ der Kranke rufen. Wie ein Ertrinkender sich an den letzten Strohhalm klammert, so hafteten des Capitäns Blicke an dem Jungen, als dieser eintrat.

„Karl Müller, hast du eine Mutter zu Hause, ist sie eine fromme Frau, hat sie dich beten gelehrt?“ „Ja, Capitän, das hat sie und sie hat mir, als ich von Hause abreiste, eine Bibel geschenkt und“ —. „Hast du die Bibel hier, Junge?“ unterbrach ihn der Capitän. „Ja, Capitän, und ich habe meiner Mutter versprechen müssen, täglich darin zu lesen.“ „Hol sie her und lies mir etwas vor, das beim Sterben hilft.“ Karl holte seine Bibel und fragte: „Was soll ich denn aufschlagen, Capitän, soll ich das lesen, was ich beim Abschied mit meiner Mutter zusammen gelesen habe?“ „Ja, ja, fang nur an.“ Da schlug Karl das 53. Capitel des Propheten Jesaias auf und las es vor; als er aber an den fünften Vers kam, stockte er. „Capitän“, fragte er, „soll ich den Vers so lesen, wie ihn mich meine Mutter lesen lehrte?“ Da las denn Karl weiter: „Aber er ist um Karl Müllers Missethat willen verwundet und um Karl Müllers Sünde willen zerschlagen. Die Strafe liegt auf ihm, auf daß Karl Müller Friede hätte, und durch seine Wunden ist Karl Müller geheilet.“ Da richtete sich der Kranke plötzlich auf: „Halt“, rief er, „das ist’s ja, was ich brauche; lies den Vers noch einmal, Junge, aber setze meinen Namen statt des deinigen.“ Da las Karl den Bibelvers noch einmal langsam und deutlich vor. Und als er geendet hatte: „Die Strafe liegt auf ihm, auf daß Peter Jensen Friede hätte, und durch seine Wunden ist Peter Jensen geheilet“, da athmete der Kranke tief auf und sprach leise vor sich hin: „Da habe ich, was ich so lange gesucht habe; das gibt Ruhe und Trost.“

Von diesem Tage an wurde Karl Müller noch manchesmal in die Kajüte des Capitäns gerufen und mußte diesem aus seiner Bibel vorlesen. Immer wieder aber wollte der Sterbende, mit dem es rasch zu Ende ging, die Worte hören: „durch seine Wunden ist Peter Jensen geheilet“. Mit diesem Troste auf den Lippen und im Herzen ging Peter noch in der elften Stunde in den himmlischen Friedenshafen ein.
